\begin{textsample}{POS Dim 1 – 2010 – Score 31.00 – 2010/t002836.txt}  \label{ex:f1_pos_026}
In \textbf{the} spirit of Jacques Cousteau, who said,"People protect what they love,"I want to share with you today what I love most in \textbf{the} ocean, and that ' s \textbf{the} incredible number and variety of animals in it that make light.
My addiction began with this strange looking diving suit called Wasp ; that ' s not an acronym — just somebody thought it looked like \textbf{the} insect.
It was actually developed for use by \textbf{the} offshore \textbf{oil} industry for diving on \textbf{oil} rigs down to a depth of 2, 000 feet.
Right after I completed my Ph.
D., I was lucky enough to be included with a group of scientists that was using it for \textbf{the} first time as a tool for ocean exploration.
We trained in a tank in Port Hueneme, and then my first open ocean dive was in Santa Barbara Channel.
It was an evening dive.
I went down to a depth of 880 feet and turned out \textbf{the} lights.
And \textbf{the} reason I turned out \textbf{the} lights is because I knew I would see this phenomenon of animals making light called bioluminescence.
But I was totally unprepared for how much there was and how spectacular it was.
I saw chains of jellyfish called siphonophores that were longer than this room, pumping out so much light that I could read \textbf{the} dials and gauges inside \textbf{the} suit without a flashlight ; and puffs and billows of what looked like luminous \textbf{blue} smoke ; and \textbf{explosions} of sparks that would swirl up out of \textbf{the} thrusters — just like when you throw a log on a campfire and \textbf{the} embers swirl up off \textbf{the} campfire, but these were icy, \textbf{blue} embers.
It was breathtaking.
Now, usually if people are familiar with bioluminescence at all, it ' s these guys ; it ' s fireflies.
And there are a few other land-dwellers that can make light — some insects, earthworms, fungi — but in general, on land, it ' s really rare.
In \textbf{the} ocean, it ' s \textbf{the} rule rather than \textbf{the} exception.
If I go out in \textbf{the} open ocean environment, virtually anywhere in \textbf{the} world, and I drag a net from 3, 000 feet to \textbf{the} \textbf{surface}, most of \textbf{the} animals — in fact, in many places, 80 to 90 percent of \textbf{the} animals that I bring up in that net — make light.
This makes for some pretty spectacular light shows.
Now I want to share with you a little video that I shot from a submersible.
I first developed this \textbf{technique} working from a little single-person submersible called Deep Rover and then adapted it for use on \textbf{the} Johnson Sea-Link, which you see here.
So, mounted in front of \textbf{the} observation \textbf{sphere}, there ' s a a three-foot \textbf{diameter} hoop with a screen stretched across it.
And inside \textbf{the} \textbf{sphere} with me is an intensified camera that ' s about as sensitive as a fully dark-adapted human eye, albeit a little fuzzy.
So you turn on \textbf{the} camera, turn out \textbf{the} lights.
That sparkle you ' re seeing is not luminescence, that ' s just electronic noise on these super intensified cameras.
You don ' t see luminescence until \textbf{the} submersible begins to move forward through \textbf{the} water, but as it does, animals bumping into \textbf{the} screen are stimulated to bioluminesce.
Now, when I was first doing this, all I was trying to do was count \textbf{the} numbers of sources.
I knew my forward speed, I knew \textbf{the} area, and so I could figure out how many hundreds of sources there were per cubic meter.
But I started to realize that I could actually identify animals by \textbf{the} type of flashes they produced.
And so, here, in \textbf{the} Gulf of Maine at 740 feet, I can name pretty much everything you ' re seeing there to \textbf{the} species level.
Like those big \textbf{explosions}, sparks, are from a little comb jelly, and there ' s krill and other kinds of crustaceans, and jellyfish.
There was another one of those comb jellies.
And so I ' ve worked with \textbf{computer} image analysis engineers to develop automatic recognition systems that can identify these animals and then extract \textbf{the} XYZ coordinate of \textbf{the} initial impact point.
And we can then do \textbf{the} kinds of things that ecologists do on land, and do nearest neighbor distances.
But you don ' t always have to go down to \textbf{the} depths of \textbf{the} ocean to see a light show like this.
You can actually see it in \textbf{surface} waters.
This is some shot, by Dr.
Mike Latz at Scripps Institution, of a dolphin swimming through bioluminescent plankton.
And this isn ' t someplace exotic like one of \textbf{the} bioluminescent \textbf{bays} in Puerto Rico, this was actually shot in San Diego Harbor.
And sometimes you can see it even closer than that, because \textbf{the} heads on ships — that ' s toilets, for any land lovers that are listening — are flushed with unfiltered seawater that often has bioluminescent plankton in it.
So, if you stagger into \textbf{the} head late at night and you ' re so toilet-hugging sick that you forget to turn on \textbf{the} light, you may think that you ' re having a religious experience.
So, how does a living creature make light?
Well, that was \textbf{the} question that 19th century French physiologist Raphael Dubois, asked about this bioluminescent clam.
He ground it up and he managed to get out a couple of chemicals ; one, \textbf{the} enzyme, he called luciferase ; \textbf{the} substrate, he called luciferin after Lucifer \textbf{the} Lightbearer.
That terminology has stuck, but it doesn ' t actually refer to specific chemicals because these chemicals come in a lot of different shapes and forms.
In fact, most of \textbf{the} people studying bioluminescence today are focused on \textbf{the} \textbf{chemistry}, because these chemicals have proved so incredibly valuable for developing antibacterial agents, cancer fighting drugs, testing for \textbf{the} presence of life on Mars, \textbf{detecting} pollutants in our waters — which is how we use it at ORCA.
In 2008, \textbf{the} Nobel Prize in \textbf{Chemistry} was awarded for work done on a molecule called green fluorescent protein that was isolated from \textbf{the} bioluminescent \textbf{chemistry} of a jellyfish, and it ' s been equated to \textbf{the} invention of \textbf{the} microscope, in terms of \textbf{the} impact that it has had on cell \textbf{biology} and genetic engineering.
Another thing all these molecules are telling us that, apparently, bioluminescence has \textbf{evolved} at least 40 times, maybe as many as 50 separate times in \textbf{evolutionary} history, which is a clear indication of how spectacularly important this trait is for survival.
So, what is it about bioluminescence that ' s so important to so many animals?
Well, for animals that are trying to avoid predators by staying in \textbf{the} darkness, light can still be very useful for \textbf{the} three basic things that animals have to do to survive : and that ' s find food, attract a \textbf{mate} and avoid being eaten.
So, for example, this \textbf{fish} has a built-in headlight behind its eye that it can use for finding food or attracting a \textbf{mate}.
And then when it ' s not using it, it actually can roll it down into its head just like \textbf{the} headlights on your Lamborghini.
This \textbf{fish} actually has high beams.
And this \textbf{fish}, which is one of my favorites, has three headlights on each side of its head.
Now, this one is \textbf{blue}, and that ' s \textbf{the} color of most bioluminescence in \textbf{the} ocean because \textbf{evolution} has selected for \textbf{the} color that travels farthest through seawater in order to optimize communication.
So, most animals make \textbf{blue} light, and most animals can only see \textbf{blue} light, but this \textbf{fish} is a really fascinating exception because it has two red light organs.
And I have no idea why there ' s two, and that ' s something I want to solve some day — but not only can it see \textbf{blue} light, but it can see red light.
So it uses its red bioluminescence like a sniper ' s scope to be able to sneak up on animals that are blind to red light and be able to see them without being seen.
It ' s also got a little chin barbel here with a \textbf{blue} luminescent lure on it that it can use to attract prey from a long way off.
And a lot of animals will use their bioluminescence as a lure.
This is another one of my favorite \textbf{fish}.
This is a viperfish, and it ' s got a lure on \textbf{the} end of a long \textbf{fishing} rod that it arches in front of \textbf{the} toothy jaw that gives \textbf{the} viperfish its name. \textbf{The} teeth on this \textbf{fish} are so long that if they closed inside \textbf{the} mouth of \textbf{the} \textbf{fish}, it would actually impale its own brain.
So instead, it \textbf{slides} in grooves on \textbf{the} outside of \textbf{the} head.
This is a Christmas tree of a \textbf{fish} ; everything on this \textbf{fish} lights up, it ' s not just that lure.
It ' s got a built-in flashlight.
It ' s got these jewel-like light organs on its belly that it uses for a type of camouflage that obliterates its shadow, so when it ' s swimming around and there ' s a predator looking up from below, it makes itself \textbf{disappear}.
It ' s got light organs in \textbf{the} mouth, it ' s got light organs in every single scale, in \textbf{the} \textbf{fins}, in a mucus layer on \textbf{the} back and \textbf{the} belly, all used for different things — some of which we know about, some of which we don ' t.
And we know a little bit more about bioluminescence thanks to Pixar, and I ' m very grateful to Pixar for sharing my favorite topic with so many people.
I do wish, with their budget, that they might have spent just a tiny bit more money to pay a consulting fee to some poor, starving graduate student, who could have told them that those are \textbf{the} eyes of a \textbf{fish} that ' s been preserved in formalin.
These are \textbf{the} eyes of a living anglerfish.
So, she ' s got a lure that she sticks out in front of this living mousetrap of needle-sharp teeth in order to attract in some unsuspecting prey.
And this one has a lure with all kinds of little interesting threads coming off it.
Now we used to think that \textbf{the} different shape of \textbf{the} lure was to attract different types of prey, but then stomach content analyses on these \textbf{fish} done by scientists, or more likely their graduate students, have revealed that they all eat pretty much \textbf{the} same thing.
So, now we believe that \textbf{the} different shape of \textbf{the} lure is how \textbf{the} male recognizes \textbf{the} female in \textbf{the} anglerfish world, because many of these males are what are known as \textbf{dwarf} males.
This little guy has no visible means of self-support.
He has no lure for attracting food and no teeth for eating it when it gets there.
His only hope for existence on this \textbf{planet} is as a gigolo.
He ' s got to find himself a babe and then he ' s got to latch on for life.
So this little guy has found himself this babe, and you will note that he ' s had \textbf{the} good sense to attach himself in a way that he doesn ' t actually have to look at her.
But he still knows a good thing when he sees it, and so he \textbf{seals} \textbf{the} relationship with an eternal kiss.
His flesh fuses with her flesh, her bloodstream grows into his body, and he becomes nothing more than a little sperm sac.
Well, this is a deep- \textbf{sea} version of Women ' s Lib.
She always knows where he is, and she doesn ' t have to be monogamous, because some of these females come up with multiple males attached.
So they can use it for finding food, for attracting \textbf{mates}.
They use it a lot for defense, many different ways.
A lot of them can release their luciferin or luferase in \textbf{the} water just \textbf{the} way a squid or an octopus will release an ink cloud.
This \textbf{shrimp} is actually spewing light out of its mouth like a fire breathing dragon in order to blind or distract this viperfish so that \textbf{the} \textbf{shrimp} can swim away into \textbf{the} darkness.
And there are a lot of different animals that can do this : There ' s jellyfish, there ' s squid, there ' s a whole lot of different crustaceans, there ' s even \textbf{fish} that can do this.
This \textbf{fish} is called \textbf{the} shining tubeshoulder because it actually has a tube on its shoulder that can squirt out light.
And I was luck enough to capture one of these when we were on a trawling expedition off \textbf{the} northwest coast of Africa for"Blue Planet,"for \textbf{the} deep portion of"Blue Planet."And we were using a special trawling net that we were able to bring these animals up alive.
So we captured one of these, and I brought it into \textbf{the} lab.
So I ' m holding it, and I ' m about to touch that tube on its shoulder, and when I do, you ' ll see bioluminescence coming out.
But to me, what ' s shocking is not just \textbf{the} amount of light, but \textbf{the} fact that it ' s not just luciferin and luciferase.
For this \textbf{fish}, it ' s actually whole cells with nuclei and membranes.
It ' s energetically very costly for this \textbf{fish} to do this, and we have no idea why it does it — another one of these great mysteries that needs to be solved.
Now, another form of defense is something called a burglar alarm — same reason you have a burglar alarm on your car ; \textbf{the} honking horn and flashing lights are meant to attract \textbf{the} attention of, hopefully, \textbf{the} police that will come and take \textbf{the} burglar away — when an animal ' s caught in \textbf{the} clutches of a predator, its only hope for escape may be to attract \textbf{the} attention of something bigger and nastier that will attack their attacker, thereby affording them a chance for escape.
This jellyfish, for example, has a spectacular bioluminescent display.
This is us chasing it in \textbf{the} submersible.
That ' s not luminescence, that ' s reflected light from \textbf{the} gonads.
We capture it in a very special \textbf{device} on \textbf{the} front of \textbf{the} submersible that allows us to bring it up in really pristine condition, bring it into \textbf{the} lab on \textbf{the} ship.
And then to generate \textbf{the} display you ' re about to see, all I did was touch it once per second on its nerve ring with a sharp pick that ' s sort of like \textbf{the} sharp tooth of a \textbf{fish}.
And once this display gets going, I ' m not touching it anymore.
This is an unbelievable light show.
It ' s this pinwheel of light, and I ' ve done calculations that show that this could be seen from as much as 300 feet away by a predator.
And I thought,"You know, that might actually make a pretty good lure."Because one of \textbf{the} things that ' s frustrated me as a deep-sea explorer is how many animals there probably are in \textbf{the} ocean that we know nothing about because of \textbf{the} way we explore \textbf{the} ocean. \textbf{The} primary way that we know about what lives in \textbf{the} ocean is we go out and drag nets behind ships.
And I defy you to name any other branch of science that still depends on hundreds of \textbf{year-old} technology. \textbf{The} other primary way is we go down with submersibles and remote-operated vehicles.
I ' ve made hundreds of dives in submersibles.
When I ' m sitting in a submersible though, I know that I ' m not unobtrusive at all — I ' ve got bright lights and noisy thrusters — any animal with any sense is going to be long gone.
So, I ' ve wanted for a long time to figure out a different way to explore.
And so, sometime ago, I got this idea for a camera system.
It ' s not exactly rocket science.
We call this thing Eye-in-the-Sea.
And scientists have done this on land for years ; we just use a color that \textbf{the} animals can ' t see and then a camera that can see that color.
You can ' t use infrared in \textbf{the} \textbf{sea}.
We use far-red light, but even that ' s a problem because it gets absorbed so quickly.
Made an intensified camera, wanted to make this electronic jellyfish.
Thing is, in science, you basically have to tell \textbf{the} funding agencies what you ' re going to discover before they ' ll give you \textbf{the} money.
And I didn ' t know what I was going to discover, so I couldn ' t get \textbf{the} funding for this.
So I kluged this together, I got \textbf{the} Harvey Mudd Engineering Clinic to actually do it as an undergraduate student project initially, and then I kluged funding from a whole bunch of different sources.
Monterey Bay Aquarium Research Institute gave me time with their ROV so that I could test it and we could figure out, you know, for example, which colors of red light we had to use so that we could see \textbf{the} animals, but they couldn ' t see us — get \textbf{the} electronic jellyfish working.
And you can see just what a shoestring operation this really was, because we cast these 16 \textbf{blue} LEDs in epoxy and you can see in \textbf{the} epoxy mold that we used, \textbf{the} word Ziploc is still visible.
Needless to say, when it ' s kluged together like this, there were a lot of trials and tribulations getting this working.
But there came a moment when it all came together, and everything worked.
And, remarkably, that moment got caught on film by photographer Mark Richards, who happened to be there at \textbf{the} precise moment that we discovered that it all came together.
That ' s me on \textbf{the} \textbf{left}, my graduate student at \textbf{the} time, Erika Raymond, and Lee Fry, who was \textbf{the} engineer on \textbf{the} project.
And we have this photograph posted in our lab in a place of honor with \textbf{the} caption :"Engineer satisfying two women at once."And we were very, very happy.
So now we had a system that we could actually take to some place that was kind of like an oasis on \textbf{the} bottom of \textbf{the} ocean that might be patrolled by large predators.
And so, \textbf{the} place that we took it to was this place called a Brine Pool, which is in \textbf{the} northern part of \textbf{the} Gulf of Mexico.
It ' s a magical place.
And I know this footage isn ' t going to look like anything to you — we had a crummy camera at \textbf{the} time — but I was ecstatic.
We ' re at \textbf{the} edge of \textbf{the} Brine Pool, there ' s a \textbf{fish} that ' s swimming towards \textbf{the} camera.
It ' s clearly undisturbed by us.
And I had my window into \textbf{the} deep \textbf{sea}.
I, for \textbf{the} first time, could see what animals were doing down there when we weren ' t down there disturbing them in some way.
Four hours into \textbf{the} deployment, we had programmed \textbf{the} electronic jellyfish to come on for \textbf{the} first time.
Eighty-six seconds after it went into its pinwheel display, we recorded this : This is a squid, over six feet long, that is so new to science, it cannot be placed in any known scientific family.
I could not have asked for a better proof of concept.
And based on this, I went back to \textbf{the} National Science Foundation and said,"This is what we will discover."And they gave me enough money to do it right, which has involved developing \textbf{the} world ' s first deep-sea webcam — which has been installed in \textbf{the} Monterey Canyon for \textbf{the} past year — and now, more recently, a modular form of this system, a much more mobile form that ' s a lot easier to launch and recover, that I hope can be used on Sylvia ' s"hope spots"to help explore and protect these areas, and, for me, learn more about \textbf{the} bioluminescence in these"hope spots."So one of these take-home messages here is, there is still a lot to explore in \textbf{the} oceans.
And Sylvia has said that we are destroying \textbf{the} oceans before we even know what ' s in them, and she ' s right.
So if you ever, ever get an opportunity to take a dive in a submersible, say yes — a thousand times, yes — and please turn out \textbf{the} lights.
I promise, you ' ll love it.
Thank you.

% matched lemmas: bay, biology, blue, chemistry, computer, detect, device, diameter, disappear, dwarf, evolution, evolutionary, evolve, explosion, fin, fish, fishing, left, mate, oil, planet, sea, seal, shrimp, slide, sphere, surface, technique, the, year-old
\end{textsample}
